\newpage
\phantomsection
\label{prf:riemann-cauchy-criterion}
\LRAProofFor{thm:riemann-cauchy-criterion}
\begin{remark*}[Return]\hyperref[thm:riemann-cauchy-criterion]{Return to Theorem}\end{remark*}
\begin{theorem*}[Cauchy Criterion for Riemann Integrability]
A bounded function is Riemann-integrable if and only if sufficiently fine
tagged sums are mutually close.
\end{theorem*}
\begin{proof}[Professional Standard Proof]
\LRAProofBodyStart
Suppose \(f\) is Riemann-integrable with value \(L\). Given \(\varepsilon>0\),
choose \(\delta>0\) so that every tagged partition of mesh less than
\(\delta\) has sum within \(\varepsilon/2\) of \(L\). Then any two such sums
differ by less than \(\varepsilon\) by the triangle inequality.

Conversely, assume the mutual closeness condition. The tagged sums, directed
by refinement and decreasing mesh, form a Cauchy net in \(\mathbb R\). Since
\(\mathbb R\) is complete, this net converges to some \(L\). The Cauchy
condition then says exactly that all sufficiently fine tagged sums are close
to this \(L\), which is the Riemann integral.
\end{proof}
\begin{proof}[Detailed Learning Proof]
\LRAProofBodyStart
If all fine sums are near \(L\), then any two are near each other because both
sit in the same small interval around \(L\). The converse uses the same idea
as the Cauchy criterion for sequences: a family whose sufficiently late values
are mutually close has a limit in a complete space.
\end{proof}
\begin{remark*}[Proof structure]
Use the triangle inequality for the easy direction and completeness for the
converse.
\end{remark*}
\begin{dependencies}
\begin{itemize}
  \item \hyperref[def:riemann-sum]{Riemann Sum}.
  \item \hyperref[def:riemann-integral]{Riemann Integral}.
\end{itemize}
\end{dependencies}
\clearpage
